% ---
% id: meat-core/frank-wolfe/req:first-order-clamp
% kind: req
% title: "Clamping of the first-order bound at zero"
% summary: >-
%   What reaches the key length is max{0, h_lb}, and the engine reports the
%   unclamped h_lb beside a flag recording the clamp. The clamp substitutes one
%   valid lower bound for another and repairs nobody's data.
% prominence: standard
% status: implemented
% uses: [thm:firstorder, cor:fo-keylength]
% mentions: [rem:fo-negative-values, rem:fo-arithmetic]
% assumes: []
% refs: []
% provenance: newly-derived
% implements: certiqs_key.solvers.frank_wolfe.general_keyrate:key_length_from_first_order
% lean:
% ---

A computed $\hlb$ may be negative, the linearisation being weak far from the
optimum or the arithmetic crossing zero, and $\max\{0,\hlb\}$ is still a
lower bound on $\hlow$ by Theorem~\ref{thm:firstorder}\,(ii),
respectively~(iiib) (Remark~\ref{rem:fo-negative-values}); the clamp corrects
only the low side (Remark~\ref{rem:fo-arithmetic}\,(a)). The unclamped value
lets a reader see what was computed and not only what was used.

\begin{requirement}\label{req:first-order-clamp}
The engine passes $\max\{0,\hlb\}$ to the key length of
Corollary~\ref{cor:fo-keylength}, and records the clamp: it emits the
boolean \iface{first_order_clamped} beside \iface{h_lb}, and \iface{h_lb}
carries the unclamped value.
\end{requirement}
